\chapter{Diskussion}
\label{ch:diskussion}

\section{Zusammenfassung und Analyse}
\label{sec:zusammenfassung}

In der \csec{A1} wurden qualitativ die Eigenschaften des Versuchsaufbaus untersucht, es zeigte sich,
eine größere Spaltbreite führt zu einem breiteren Objektbild und entsprechend zu einem schmaleren Beugungsbild. Zudem bewirkt eine Rotation des Spalts eine gleichartige Rotation sowohl des Objekt- als auch des Beugungsbildes.
Eine breitere Öffnung des Analysespalts führt bei kritischer Spaltbreite dazu, dass im Objektbild zwei Lichtlinien sichtbar werden.

\subsection*{Einzelspalt}

In der \csec{A2} wurde eine Bestimmung der Extremapositionen sowie deren relativer Intensitäten geliefert und Vergleich mit der Theorie gestellt. 
Zunächst wurde die Eichgerade anhand von \cabb{plots/Eichgerade2a} bestimmt. Daraus ergibt sich:
\imp[-6]{s_\text{E}}{2.59}{0.03}{\frac{m}{px}}
Die Abstände der Extrema in Abhängigkeit von ihrer Ordnung wurden in \cabb{plots/Extrema_Abstand_Ordnungszahltry2} dargestellt. Die Steigung ergibt sich zu:
\imp{s_\text{ord}}{4.71}{0.07}{\cdot 10 , px}
Die Bestimmung der experimentellen Ordnung zeigt laut \ctab{vergleich_ordnungszahl_theo_ex_einzelspalt} eine gute Übereinstimmung mit der Theorie; alle Werte liegen innerhalb von $3\sigma$.
Die aus den beiden Steigungen sowie der Wellenlänge des Lichts und der Brennweite der Linse $L_1$ bestimmte Spaltbreite beträgt:
\res[-5]{b}{6.96}{0.14}{m}
Anschließend wurden die relativen Intensitäten, bezogen auf das Hauptmaximum, bestimmt. Die Ergebnisse aus \ctab{intensitaeten} zeigen jedoch keine gute Übereinstimmung mit der Theorie; alle Werte weichen signifikant ab. Auffällig ist, dass sämtliche Werte systematisch zu klein sind und die relativen Abweichungen vergleichbare Größenordnungen aufweisen. Dies deutet auf einen systematischen Fehler hin, beispielsweise einen unpräzisen Versuchsaufbau oder eine fehlerhafte Einstellung der Messsoftware.

In \csec{A4} wurde das Objektbild als Fouriersynthese des Beugungsbildes betrachtet. Zur Bestimmung der Spaltbreite wurden Bild- und Gegenstandsweite getrennt ausgewertet, die theoretisch denselben Wert liefern sollten. Es ergibt sich jedoch über die Bildweite $b$:
\res[-4]{d_{1,b}}{3.68}{0.07}{m}
und über die Gegenstandsweite:
\res[-4]{d_{1,g}}{5.44}{0.09}{m}

Die Abweichung beträgt $15.00\sigma$ und ist damit hochsignifikant. Dies deutet erneut auf einen unpräzisen Versuchsaufbau hin.

In der folgenden Aufgabe wurden Maxima nur bis zu bestimmten Ordnungen zugelassen: bis zur 0., 1., 2. und fälschlicherweise bis zur 4. Ordnung (korrekt wäre die 3. Ordnung gewesen).
Der Vergleich der theoretischen und experimentellen Kurven in \cabb{plots/vergleich_maxima_3a} zeigt eine grundsätzliche Übereinstimmung der Verläufe. Allerdings sind die experimentellen Kurven nach links verschoben, was auf eine nicht mittig platzierte Kamera hindeutet. Zudem zeigt sich eine Asymmetrie, die auf eine Verkippung des Spalts schließen lässt.
Die relativen Intensitäten wurden erneut bestimmt. Dabei zeigt sich, dass die Intensitäten bis zur 1. Ordnung signifikant von der Theorie abweichen, während sie für die 2. Ordnung innerhalb des nicht signifikanten Bereichs liegen. Dieser Trend entspricht der theoretischen Erwartung, da mit zunehmender Ordnung die Übereinstimmung besser wird. Allerdings wird dieser Verlauf durch die (fälschlich betrachtete) 4. Ordnung unterbrochen, bei der erneut signifikante Abweichungen auftreten. Dies spricht für eine fehlerhafte Einstellung der Messsoftware, insbesondere eine unpassende Belichtungszeit.

\subsection*{Doppelspalt}

In \csec{A3} wurde der Doppelspalt untersucht. Zunächst erfolgte die Vermessung der Spaltbreite über Bild- und Gegenstandsweite. Über die Bildweite ergibt sich:
\res[-8]{b_\text{b,dop}}{2.57}{0.06}{m}
Über die Gegenstandsweite ergibt sich hingegen:
\res[-8]{b_\text{g,dop}}{3.79}{0.08}{m}
Die Abweichung beträgt gemäß \hyperref[eq:standardabweichung]{Standardabweichung (\ref*{eq:standardabweichung})} $11.93\sigma$ und ist damit signifikant.

Zusätzlich wurde der Abstand zwischen den Spalten bestimmt. Über die Bildweite ergibt sich:
\res[-8]{d_\text{b,dop}}{5.99}{0.12}{m}
und über die Gegenstandsweite:
\res[-8]{d_\text{g,dop}}{8.84}{0.14}{m}

Die Abweichung beträgt hier $14.93\sigma$ und ist ebenfalls signifikant. Auch dies weist auf einen nicht korrekt justierten Versuchsaufbau hin, da die aus Objekt- und Bildweite bestimmten Werte deutlich voneinander abweichen.
Für den Doppelspalt wurden ebenfalls relative Intensitäten bestimmt, sowohl mittels Python-Auswertung als auch durch Cursor-Messung. Die Ergebnisse sind in \ctab{doppelspalt_int} dargestellt. Es zeigt sich erneut eine Asymmetrie; lediglich das erste Maximum weist keine signifikante Abweichung auf. Insgesamt kann die Theorie anhand dieser Messwerte kaum bestätigt werden.

Abschließend wurde in \csec{A5} das Objektbild als Fouriersynthese des Beugungsbildes für den Doppelspalt untersucht. Dabei wurden Maxima bis zur 0., 1. und 2. Ordnung berücksichtigt. Für die nullte Ordnung ergibt sich die theoretische kritische Ordnung zu:
\[
    n_\text{krit} = (0.256 \pm 0.008)
\]
Für die kritische sowie die erste Ordnung konnten Grenzfrequenzen bestimmt werden, die in \ctab{finalllllllalllllalllllallllll} aufgeführt sind. Es zeigt sich, dass die experimentellen Werte nicht mit den theoretischen übereinstimmen, da alle Abweichungen statistisch signifikant sind.

\section{Kritik}

Der Versuchsaufbau wurde insgesamt unzureichend umgesetzt. Insbesondere war die Dokumentation im Messprotokoll lückenhaft und unpräzise. Wichtige Einstellungen der Messsoftware, wie beispielsweise die Belichtungszeit, hätten sorgfältig dokumentiert werden müssen.
